\section{Failure of separate variance and form comparison}\label{separatecomparison}
The regularized spectrum does not supply a uniform comparison of the
stationary variances. The obstruction below concerns actual chains,
including their nonlinear reconstructed potentials. For an energy-normalized
positive potential $v$, write
\[
 \mathcal E_v(f)=\frac12\E_{x\sim\mu}\sum_y
 P(x,y)v(x)v(y)(f(x)-f(y))^2.
\]
This is the Dirichlet form of $K_v$ in stationary measure $vPv\,d\mu$.

\begin{theorem}[Separate-comparison obstruction]\label{separatecomparisonfailure}
Fix $0<\rho,b<1$, $C>0$, and $0\leq\tau<1$. There are positive
energy-normalized potentials $h_L,u_L$, with
\[
 \pi_L=h_LPh_L,\qquad w_L=u_LPu_L=P_\tau\pi_L,
\]
both satisfying \eqref{admissible}, and events $A_L$, such that, for
$f_L=1_{A_L}$,
\[
 \frac{\operatorname{Var}_{\pi_L}(f_L)}{\operatorname{Var}_{w_L}(f_L)}
 \longrightarrow\infty,\qquad
 \frac{\mathcal E_{h_L}(f_L)}{\mathcal E_{u_L}(f_L)}
 \longrightarrow\infty.
\]
Nevertheless both indicator Rayleigh quotients tend to one, and both
full-chain gaps tend to $1-\rho$. The family can be chosen with
$\norm{\pi_L-w_L}_s\to c>0$ and with $\pi_L^s$ not uniformly
integrable, despite $\E\pi_L^s<2$.
\end{theorem}
\begin{proof}
Put $B=2^{1/s}$ and $c=\min\{(1-b)/2,(B-1)/4\}>0$.
For all sufficiently large $L$, set
\[
 d=k_L+1,\qquad n=\lfloor L/d\rfloor,\qquad p=2^{-n},
 \qquad t=cp^{-1/s}.
\]
Take $n$ disjoint blocks of size $d$, let $z_i$ be their parities,
and set $A=\{z_1=\cdots=z_n=1\}$. Remaining coordinates are unused.
Thus $\mu(A)=p$, $d\sim C\sqrt L$, and $n\sim\sqrt L/C$.
Define
\[
 \pi=1+t(1_A-p).
\]
It has mean one and $\pi\geq1-c>b$. Moreover
\[
 \norm\pi_s\leq1+t(p^{1/s}+p)
       =1+c(1+p^{1-1/s})\leq1+2c<B.
\]
The expansion $1_A=p\prod_{i=1}^n(1+z_i)$ has nonconstant Fourier
support only on unions of complete blocks, whose weights are at least
$d>k_L$. This proves every labelled matching condition. The density
$w=P_\tau\pi$ has the same mean, stationary floor, cap and matching,
by positivity, $L^s$ contraction and the diagonal Fourier action.
Lemma~\ref{inversestationary} gives the unique positive solutions
$hPh=\pi$ and $uPu=w$, automatically of energy one. In the finitely
many excluded dimensions take $h=u=1$ and any nontrivial event.

The exact smoothed density is
\[
 w(z)=1+tp\left[\prod_{i=1}^n(1+\tau^d z_i)-1\right].
\]
For $0\leq\epsilon<1$ the maximum absolute deviation of this product
from one is $(1+\epsilon)^n-1$: its two extremes are
$(1\pm\epsilon)^n$, whose sum is at least two. Hence
\[
 \norm{w-1}_\infty
 \leq tp\big[(1+\tau^d)^n-1\big]
 \leq cp^{1-1/s}\big[\exp(n\tau^d)-1\big]\longrightarrow0.
\]
Here $n\tau^d\to0$ for each fixed $\tau<1$, including $\tau=0$.
No assertion uniform in $\tau$ approaching one is required.

Both reconstructed potentials belong to the same block-parity algebra.
Indeed translations having parity one on each block, and arbitrary
translations on unused coordinates, leave $\pi$ and $w$ invariant and
commute with $P$. Translating a positive reconstructed potential therefore
gives another solution of the same equation. Uniqueness in
Lemma~\ref{inversestationary} proves invariance. The translation orbits
are exactly the parity fibers, so $h=H(z)$ and $u=J(z)$.

We record actual return probabilities for either reconstructed potential.
Write $H$ for any positive block-parity potential of energy one and
$G=PH$. Put
\[
 \theta=\rho^d,\quad \ell=(1-\theta)^n,\quad v=(1+\theta)^n,
 \quad a=\E H.
\]
The quotient noise transition from $z$ to $z'$ is
$p\prod_i(1+\theta z_i z'_i)$. Its density relative to uniform
quotient measure lies between $\ell$ and $v$. Consequently
\[
 a\ell\leq G\leq av,\qquad
 v^{-1}\leq a^2\leq\ell^{-1},\qquad
 \frac{\ell^2}{v}\leq G^2\leq\frac{v^2}{\ell}.
\]
The middle bounds use $1=\E HG$. If $q$ is the stationary mass of
$A$, then $q=pH(+)G(+)$. The return probability from any point of
$A$ to the \emph{whole fiber} $A$ is
\[
 r_A=K_H(x,A)=\frac{pvH(+)}{G(+)}
              =\frac{vq}{G(+)^2}\leq\frac{v^2}{\ell^2}q.
\]
This is not a single-vertex holding probability. Since $n\rho^d\to0$,
both $\ell$ and $v$ tend to one. For the two chains the stationary
event masses are
\[
 q_\pi=p[1+t(1-p)]\sim tp\longrightarrow0,\qquad
 q_w=p\left[1+tp\big((1+\tau^d)^n-1\big)\right]\sim p.
\]
Thus $q_\pi/q_w\sim t\to\infty$ and both actual return
probabilities $r_\pi,r_w$ tend to zero. Reversibility gives exactly
\[
 \operatorname{Var}_{\pi}(1_A)=q_\pi(1-q_\pi),\quad
 \mathcal E_h(1_A)=q_\pi(1-r_\pi),
\]
and the analogous identities with $w,u$. Both displayed ratios in the
theorem are therefore asymptotic to $t=c2^{n/s}$, whereas each
Rayleigh quotient tends to one.

To control the \emph{full} gaps, apply Proposition~\ref{blocktheorem},
which includes all nonconstant ambient sectors and unused coordinates,
to $H$ and $J$. In its notation the lower and upper bounds are
\[
 \max\{0,\rho-\varepsilon_d\}\leq\lambda_2
 \leq\min\{1,\max\{1-\alpha,\beta\}\},
\]
where, eventually $d\geq2$,
\[
 \alpha=\left(\frac{1-\rho^d}{1+\rho^d}\right)^n,\quad
 \beta=\frac{\rho+\rho^{d-1}}{1-\rho^d},\quad
 \varepsilon_d=\frac{(1-\rho^2)\rho^{d-1}}{1-\rho^d}.
\]
These tend to $1,\rho,0$, respectively. No upper bound on the
heights of $H$ or $J$ enters this proposition. Hence both full gaps
tend to $1-\rho$. This conclusion for this particular family and
every fixed $\tau<1$ does not enlarge the universal smoothing
interval in Theorem~\ref{densityspectrumtheorem}.

Finally, $\norm{t1_A}_s=c$ and
\[
 \norm{tP_\tau1_A}_s\leq tp(1+\tau^d)^n\longrightarrow0,
\]
so the reverse triangle inequality gives $\norm{\pi-w}_s\to c$.
Directly,
\[
 \E\pi^s=(1-p)(1-tp)^s+p[1+t(1-p)]^s\longrightarrow1+c^s.
\]
For every fixed threshold $R$, the value of $\pi^s$ on $A$ eventually
exceeds $R$, while its contribution to this expectation tends to $c^s$.
Thus $\pi_L^s$ is not uniformly integrable (interpreting the random
variables through their laws on the varying cubes).
\end{proof}

For clarity, a usual separate-comparison route to a gap transfer would
require both $\operatorname{Var}_\pi(f)\leq V\operatorname{Var}_w(f)$
and $\mathcal E_u(f)\leq D\mathcal E_h(f)$ uniformly.
Theorem~\ref{separatecomparisonfailure} disproves its variance premise.
It also disproves an upper bound $\mathcal E_h\leq D'\mathcal E_u$,
and hence two-sided form equivalence. It does \emph{not} disprove
$\mathcal E_u\leq D\mathcal E_h$, a direct Rayleigh-quotient
comparison, or the original gap conjecture. In particular it is not
a counterexample to comparison results whose hypotheses include a common
invariant measure, such as the setting described in \cite{alpw2022}.
The useful remaining route must retain the joint energy and variance
information, rather than impose the disproved separate variance bound.
